\documentclass{article}
\usepackage{amsmath, amssymb, geometry, enumitem, xcolor}
\geometry{a4paper, margin=1in}

\title{Math Problem Sets: Introduction to Linear Transformations}
\author{Based on Section 1.8 Concepts}
\date{}

\begin{document}

\maketitle
\tableofcontents
\newpage

% ---------------------------------------------------------
% CONCEPT CHECKLIST & QUESTION MAPPING (START)
% ---------------------------------------------------------
\section{Concept Checklist \& Question Mapping}
This checklist outlines all the core concepts and problem types extracted from the provided text, mapped to the questions in the problem set. The weight of questions is distributed with harder/foundational concepts (like Inverse Mappings and Abstract Properties) receiving more problems.

\begin{itemize}
    \item \textbf{Type 1: Forward Mapping (Image Calculation)}
    \begin{itemize}
        \item Computing $T(\mathbf{x}) = A\mathbf{x}$ for numerical vectors. [Q1, Q2]
        \item Computing $T(\mathbf{x}) = A\mathbf{x}$ for abstract/parametric vectors. [Q3, Q4]
    \end{itemize}
    
    \item \textbf{Type 2: Inverse Mapping (Finding Pre-Images \& Uniqueness)}
    \begin{itemize}
        \item Solving $A\mathbf{x} = \mathbf{b}$ for a unique solution. [Q5, Q6]
        \item Solving $A\mathbf{x} = \mathbf{b}$ with infinite solutions (parametric form). [Q7, Q9]
        \item Solving $A\mathbf{x} = \mathbf{b}$ with no solution (inconsistent). [Q8]
        \item Diagnostic: Conflating existence with uniqueness (Flawed Problem concept). [Q10]
    \end{itemize}
    
    \item \textbf{Type 3: Kernel / Mapping to the Zero Vector}
    \begin{itemize}
        \item Solving the homogeneous equation $A\mathbf{x} = \mathbf{0}$. [Q11, Q12]
        \item Parametric vector decomposition with free variables. [Q13, Q14]
    \end{itemize}
    
    \item \textbf{Type 4: Domain and Codomain Identification}
    \begin{itemize}
        \item Matrix dimensions $m \times n$ mapping $\mathbb{R}^n \to \mathbb{R}^m$. [Q16, Q17]
        \item Reversing the rule to find matrix dimensions from mapping spaces. [Q18]
        \item Diagnostic: Dimension mismatch flaw. [Q15]
    \end{itemize}
    
    \item \textbf{Type 5: Abstract Property Application (Linearity)}
    \begin{itemize}
        \item Applying Superposition: $T(c_1\mathbf{u} + c_2\mathbf{v}) = c_1T(\mathbf{u}) + c_2T(\mathbf{v})$. [Q19, Q20, Q21]
        \item Finding images of scaled and summed vectors without matrices. [Q22]
    \end{itemize}
    
    \item \textbf{Type 6: Proving Non-Linearity}
    \begin{itemize}
        \item The "Plus a Constant" Trap / Origin Rule $T(\mathbf{0}) \neq \mathbf{0}$. [Q23, Q24]
        \item Counterexamples for scalar multiplication (absolute values, powers). [Q25, Q26]
    \end{itemize}
    
    \item \textbf{Type 7: Constructing the Standard Matrix}
    \begin{itemize}
        \item Assembling $A$ using the images of standard basis vectors $T(\mathbf{e}_1), T(\mathbf{e}_2)$. [Q27, Q28]
    \end{itemize}
    
    \item \textbf{Type 8: Real-World Applications \& Modeling}
    \begin{itemize}
        \item Markov Chains and State Vectors (Finance focus). [Q29, Q30]
    \end{itemize}
\end{itemize}

\newpage

% ---------------------------------------------------------
% PROBLEM SETS
% ---------------------------------------------------------
\section{Problem Sets}

\subsection*{Type 1: Forward Mapping (Image Calculation)}
\textbf{1.} Let $A = \begin{pmatrix} 3 & -1 \\ 2 & 4 \end{pmatrix}$. Define $T: \mathbb{R}^2 \to \mathbb{R}^2$ by $T(\mathbf{x}) = A\mathbf{x}$. Find the image of $\mathbf{u} = \begin{pmatrix} 2 \\ -3 \end{pmatrix}$.

\vspace{0.3cm}
\textbf{2.} Let $A = \begin{pmatrix} 1 & 0 & -2 \\ 2 & -1 & 3 \\ 0 & 4 & 1 \end{pmatrix}$. Find the image of the vector $\mathbf{x} = \begin{pmatrix} -1 \\ 2 \\ 1 \end{pmatrix}$ under the mapping $\mathbf{x} \mapsto A\mathbf{x}$.

\vspace{0.3cm}
\textbf{3.} Let $A = \begin{pmatrix} 2 & -5 \\ 1 & 3 \\ 0 & 1 \end{pmatrix}$. Find the image of the abstract vector $\mathbf{v} = \begin{pmatrix} x_1 \\ x_2 \end{pmatrix}$.

\vspace{0.3cm}
\textbf{4.} Evaluate the matrix transformation $T(\mathbf{x}) = A\mathbf{x}$ where $A = \begin{pmatrix} 1 & 2 & 0 \\ -1 & 1 & 1 \end{pmatrix}$ and $\mathbf{x} = \begin{pmatrix} a \\ b \\ c \end{pmatrix}$.

\subsection*{Type 2: Inverse Mapping (Finding Pre-Images \& Uniqueness)}
\textbf{5.} Let $A = \begin{pmatrix} 4 & 1 \\ 2 & -3 \end{pmatrix}$ and $\mathbf{b} = \begin{pmatrix} 9 \\ -6 \end{pmatrix}$. Find a vector $\mathbf{x}$ whose image under $T(\mathbf{x}) = A\mathbf{x}$ is $\mathbf{b}$. Is this vector unique?

\vspace{0.3cm}
\textbf{6.} Let $A = \begin{pmatrix} 1 & 0 & 2 \\ 0 & 1 & -1 \\ 3 & 1 & 0 \end{pmatrix}$ and $\mathbf{b} = \begin{pmatrix} 7 \\ -2 \\ 9 \end{pmatrix}$. Find $\mathbf{x}$ such that $T(\mathbf{x}) = \mathbf{b}$.

\vspace{0.3cm}
\textbf{7.} Let $A = \begin{pmatrix} 1 & -3 & 2 \\ -2 & 6 & -4 \end{pmatrix}$ and $\mathbf{b} = \begin{pmatrix} 5 \\ -10 \end{pmatrix}$. Find the general solution for the pre-image of $\mathbf{b}$ in parametric vector form.

\vspace{0.3cm}
\textbf{8.} Let $A = \begin{pmatrix} 1 & 2 \\ 3 & 6 \\ 2 & 4 \end{pmatrix}$ and $\mathbf{b} = \begin{pmatrix} 4 \\ 12 \\ 7 \end{pmatrix}$. Does there exist a pre-image $\mathbf{x}$ such that $A\mathbf{x} = \mathbf{b}$? Explain.

\vspace{0.3cm}
\textbf{9.} Let $A = \begin{pmatrix} 1 & -2 & 1 & 0 \\ 0 & 1 & -1 & 2 \end{pmatrix}$ and $\mathbf{b} = \begin{pmatrix} 3 \\ 4 \end{pmatrix}$. Express the pre-image $\mathbf{x}$ in parametric vector form.

\vspace{0.3cm}
\textbf{10. (Find the Flaw)} Given $A = \begin{pmatrix} 1 & 3 \\ 2 & 6 \end{pmatrix}$ and $\mathbf{b} = \begin{pmatrix} 4 \\ 8 \end{pmatrix}$, a student notes that $\mathbf{x} = \begin{pmatrix} 1 \\ 1 \end{pmatrix}$ maps exactly to $\mathbf{b}$. The student writes: "Because I found a valid vector $\mathbf{x}$, the pre-image exists and is unique." Identify the specific error in this reasoning.

\subsection*{Type 3: Kernel / Mapping to the Zero Vector}
\textbf{11.} Let $A = \begin{pmatrix} 3 & -6 \\ -2 & 4 \end{pmatrix}$. Find all $\mathbf{x}$ in $\mathbb{R}^2$ that are mapped into the zero vector. 

\vspace{0.3cm}
\textbf{12.} Let $A = \begin{pmatrix} 1 & 2 & -1 \\ 2 & 4 & -2 \\ -1 & -2 & 1 \end{pmatrix}$. Find all $\mathbf{x}$ in $\mathbb{R}^3$ such that $A\mathbf{x} = \mathbf{0}$, expressed in parametric vector form.

\vspace{0.3cm}
\textbf{13.} Let $A = \begin{pmatrix} 1 & 0 & 3 & -2 \\ 0 & 1 & -4 & 1 \\ 0 & 0 & 0 & 0 \end{pmatrix}$. Determine the vectors that map to the zero vector.

\vspace{0.3cm}
\textbf{14.} Consider the transformation $T(\mathbf{x}) = A\mathbf{x}$ where $A = \begin{pmatrix} 1 & 1 & 1 & 1 \\ 2 & 2 & 2 & 2 \end{pmatrix}$. Find the parametric vector decomposition of the solution to $T(\mathbf{x}) = \mathbf{0}$. 

\subsection*{Type 4: Domain and Codomain Identification}
\textbf{15. (Find the Flaw)} A student claims: "If the transformation matrix $A$ has 5 rows and 3 columns, then $A$ is a $5 \times 3$ matrix, which means it multiplies vectors with 5 entries. Therefore, the domain is $\mathbb{R}^5$ and the codomain is $\mathbb{R}^3$." Identify the error and state the correct domain and codomain.

\vspace{0.3cm}
\textbf{16.} Let $A$ be a $7 \times 4$ matrix. What must $a$ and $b$ be in order to define $T: \mathbb{R}^a \to \mathbb{R}^b$ by $T(\mathbf{x}) = A\mathbf{x}$?

\vspace{0.3cm}
\textbf{17.} Suppose a linear transformation maps vectors from $\mathbb{R}^6$ to $\mathbb{R}^2$. What are the dimensions of its standard matrix $A$?

\vspace{0.3cm}
\textbf{18.} Define $T(\mathbf{x}) = \begin{pmatrix} x_1 - x_2 \\ x_2 \\ x_1 + x_2 \end{pmatrix}$. Identify the domain and codomain of $T$.

\subsection*{Type 5: Abstract Property Application (Linearity)}
\textbf{19.} Let $T: \mathbb{R}^2 \to \mathbb{R}^2$ be a linear transformation. If $T(\mathbf{u}) = \begin{pmatrix} 2 \\ -1 \end{pmatrix}$ and $T(\mathbf{v}) = \begin{pmatrix} 0 \\ 3 \end{pmatrix}$, find $T(3\mathbf{u} - 2\mathbf{v})$.

\vspace{0.3cm}
\textbf{20.} Let $T$ be a linear transformation where $T(\mathbf{e}_1) = \begin{pmatrix} 1 \\ 4 \end{pmatrix}$, $T(\mathbf{e}_2) = \begin{pmatrix} -2 \\ 1 \end{pmatrix}$, and $T(\mathbf{e}_3) = \begin{pmatrix} 0 \\ 5 \end{pmatrix}$. Find $T(2\mathbf{e}_1 + 3\mathbf{e}_2 - \mathbf{e}_3)$.

\vspace{0.3cm}
\textbf{21.} Given that $T(\mathbf{x}) = A\mathbf{x}$, use matrix properties to algebraically prove that $T(c_1\mathbf{v}_1 + c_2\mathbf{v}_2) = c_1T(\mathbf{v}_1) + c_2T(\mathbf{v}_2)$.

\vspace{0.3cm}
\textbf{22.} Suppose $T(\mathbf{x}_1, \mathbf{x}_2) = (3x_1 - x_2, x_1 + 4x_2)$. Verify the superposition principle $T(\mathbf{u} + \mathbf{v}) = T(\mathbf{u}) + T(\mathbf{v})$ specifically for $\mathbf{u} = (1, 1)$ and $\mathbf{v} = (2, 0)$.

\subsection*{Type 6: Proving Non-Linearity}
\textbf{23.} The "Plus a Constant" Trap: Show that the affine transformation $T(x_1, x_2) = (x_1 + 2, x_2 - 4)$ is not a linear transformation using the fastest method.

\vspace{0.3cm}
\textbf{24.} Show that the transformation $T(x_1, x_2) = (x_1^2, x_2)$ is not linear by finding a counterexample for scalar multiplication.

\vspace{0.3cm}
\textbf{25.} Show that the transformation $T(x_1, x_2) = (x_1 - x_2, |x_1|)$ is not linear. (Hint: Consider negative scalars).

\vspace{0.3cm}
\textbf{26.} Let $T(x_1, x_2) = (x_1 x_2, 3x_2)$. Prove that $T$ is not a linear transformation by showing it fails the preservation of addition property.

\subsection*{Type 7: Constructing the Standard Matrix}
\textbf{27.} Let $T: \mathbb{R}^2 \to \mathbb{R}^2$ be a linear transformation that maps $\mathbf{e}_1 = \begin{pmatrix} 1 \\ 0 \end{pmatrix}$ to $\begin{pmatrix} 4 \\ -1 \end{pmatrix}$ and maps $\mathbf{e}_2 = \begin{pmatrix} 0 \\ 1 \end{pmatrix}$ to $\begin{pmatrix} 2 \\ 7 \end{pmatrix}$. Construct the standard matrix $A$ for this transformation.

\vspace{0.3cm}
\textbf{28.} Let $T: \mathbb{R}^3 \to \mathbb{R}^4$ be a mapping where $T(\mathbf{e}_1) = \begin{pmatrix} 1 \\ 0 \\ 2 \\ -1 \end{pmatrix}$, $T(\mathbf{e}_2) = \begin{pmatrix} 3 \\ 1 \\ 0 \\ 0 \end{pmatrix}$, and $T(\mathbf{e}_3) = \begin{pmatrix} -1 \\ -1 \\ 1 \\ 1 \end{pmatrix}$. Find the matrix $A$ such that $T(\mathbf{x}) = A\mathbf{x}$.

\subsection*{Type 8: Real-World Applications \& Modeling}
\textbf{29.} Credit Rating Markov Chain: An investment bank has a transition matrix $A = \begin{pmatrix} 0.85 & 0.20 \\ 0.15 & 0.80 \end{pmatrix}$ representing the probability of a bond moving between Investment Grade (Row 1) and Speculative Grade (Row 2). If the initial portfolio is $100\%$ Investment Grade, so $\mathbf{x}_0 = \begin{pmatrix} 1 \\ 0 \end{pmatrix}$, evaluate the linear transformation $T(\mathbf{x}_0) = A\mathbf{x}_0$ to find the distribution after one year.

\vspace{0.3cm}
\textbf{30.} Three-State Markov Chain: A system moves between states A, B, and C with the transition matrix $A = \begin{pmatrix} 0.5 & 0.3 & 0.1 \\ 0.3 & 0.5 & 0.2 \\ 0.2 & 0.2 & 0.7 \end{pmatrix}$. If the current state vector is $\mathbf{x}_0 = \begin{pmatrix} 0.4 \\ 0.4 \\ 0.2 \end{pmatrix}$, calculate the expected state distribution $\mathbf{x}_1 = A\mathbf{x}_0$. Check that the sum of the probabilities in $\mathbf{x}_1$ equals $1.0$.

\newpage

% ---------------------------------------------------------
% DETAILED SOLUTIONS
% ---------------------------------------------------------
\section{Detailed Solutions}

\textbf{Solution 1:}
To find the image, compute $A\mathbf{u}$:
$$T(\mathbf{u}) = \begin{pmatrix} 3 & -1 \\ 2 & 4 \end{pmatrix} \begin{pmatrix} 2 \\ -3 \end{pmatrix} = \begin{pmatrix} (3)(2) + (-1)(-3) \\ (2)(2) + (4)(-3) \end{pmatrix} = \begin{pmatrix} 6 + 3 \\ 4 - 12 \end{pmatrix} = \begin{pmatrix} 9 \\ -8 \end{pmatrix}.$$

\textbf{Solution 2:}
Compute $A\mathbf{x}$:
$$T(\mathbf{x}) = \begin{pmatrix} 1 & 0 & -2 \\ 2 & -1 & 3 \\ 0 & 4 & 1 \end{pmatrix} \begin{pmatrix} -1 \\ 2 \\ 1 \end{pmatrix} = \begin{pmatrix} (1)(-1) + (0)(2) + (-2)(1) \\ (2)(-1) + (-1)(2) + (3)(1) \\ (0)(-1) + (4)(2) + (1)(1) \end{pmatrix} = \begin{pmatrix} -1 - 2 \\ -2 - 2 + 3 \\ 0 + 8 + 1 \end{pmatrix} = \begin{pmatrix} -3 \\ -1 \\ 9 \end{pmatrix}.$$

\textbf{Solution 3:}
Compute $A\mathbf{v}$:
$$T(\mathbf{v}) = \begin{pmatrix} 2 & -5 \\ 1 & 3 \\ 0 & 1 \end{pmatrix} \begin{pmatrix} x_1 \\ x_2 \end{pmatrix} = \begin{pmatrix} 2x_1 - 5x_2 \\ x_1 + 3x_2 \\ 0x_1 + x_2 \end{pmatrix} = \begin{pmatrix} 2x_1 - 5x_2 \\ x_1 + 3x_2 \\ x_2 \end{pmatrix}.$$

\textbf{Solution 4:}
Compute $A\mathbf{x}$:
$$T(\mathbf{x}) = \begin{pmatrix} 1 & 2 & 0 \\ -1 & 1 & 1 \end{pmatrix} \begin{pmatrix} a \\ b \\ c \end{pmatrix} = \begin{pmatrix} 1a + 2b + 0c \\ -1a + 1b + 1c \end{pmatrix} = \begin{pmatrix} a + 2b \\ -a + b + c \end{pmatrix}.$$

\textbf{Solution 5:}
We solve $A\mathbf{x} = \mathbf{b}$ via augmented matrix $[A | \mathbf{b}]$:
$$\begin{pmatrix} 4 & 1 & | & 9 \\ 2 & -3 & | & -6 \end{pmatrix}$$
$R_1 \leftrightarrow R_2$: $\begin{pmatrix} 2 & -3 & | & -6 \\ 4 & 1 & | & 9 \end{pmatrix}$ \\
$R_2 \leftarrow R_2 - 2R_1$: $\begin{pmatrix} 2 & -3 & | & -6 \\ 0 & 7 & | & 21 \end{pmatrix}$ \\
$7x_2 = 21 \implies x_2 = 3$. \\
$2x_1 - 3(3) = -6 \implies 2x_1 - 9 = -6 \implies 2x_1 = 3 \implies x_1 = 1.5$. \\
The vector is $\mathbf{x} = \begin{pmatrix} 1.5 \\ 3 \end{pmatrix}$. Since there is a pivot in every column, there are no free variables, so the solution is \textbf{unique}.

\textbf{Solution 6:}
Set up $[A | \mathbf{b}]$:
$$\begin{pmatrix} 1 & 0 & 2 & | & 7 \\ 0 & 1 & -1 & | & -2 \\ 3 & 1 & 0 & | & 9 \end{pmatrix}$$
$R_3 \leftarrow R_3 - 3R_1$: $\begin{pmatrix} 1 & 0 & 2 & | & 7 \\ 0 & 1 & -1 & | & -2 \\ 0 & 1 & -6 & | & -12 \end{pmatrix}$ \\
$R_3 \leftarrow R_3 - R_2$: $\begin{pmatrix} 1 & 0 & 2 & | & 7 \\ 0 & 1 & -1 & | & -2 \\ 0 & 0 & -5 & | & -10 \end{pmatrix}$ \\
$-5x_3 = -10 \implies x_3 = 2$. \\
$x_2 - 2 = -2 \implies x_2 = 0$. \\
$x_1 + 2(2) = 7 \implies x_1 = 3$. \\
Answer: $\mathbf{x} = \begin{pmatrix} 3 \\ 0 \\ 2 \end{pmatrix}$.

\textbf{Solution 7:}
Set up $[A | \mathbf{b}]$:
$$\begin{pmatrix} 1 & -3 & 2 & | & 5 \\ -2 & 6 & -4 & | & -10 \end{pmatrix}$$
$R_2 \leftarrow R_2 + 2R_1$: $\begin{pmatrix} 1 & -3 & 2 & | & 5 \\ 0 & 0 & 0 & | & 0 \end{pmatrix}$ \\
Basic variable: $x_1$. Free variables: $x_2, x_3$. \\
$x_1 - 3x_2 + 2x_3 = 5 \implies x_1 = 5 + 3x_2 - 2x_3$. \\
Parametric vector form: 
$$\mathbf{x} = \begin{pmatrix} 5 \\ 0 \\ 0 \end{pmatrix} + x_2 \begin{pmatrix} 3 \\ 1 \\ 0 \end{pmatrix} + x_3 \begin{pmatrix} -2 \\ 0 \\ 1 \end{pmatrix}$$

\textbf{Solution 8:}
Set up $[A | \mathbf{b}]$:
$$\begin{pmatrix} 1 & 2 & | & 4 \\ 3 & 6 & | & 12 \\ 2 & 4 & | & 7 \end{pmatrix}$$
$R_2 \leftarrow R_2 - 3R_1$ and $R_3 \leftarrow R_3 - 2R_1$:
$$\begin{pmatrix} 1 & 2 & | & 4 \\ 0 & 0 & | & 0 \\ 0 & 0 & | & -1 \end{pmatrix}$$
The last row implies $0x_1 + 0x_2 = -1$, which is impossible. Thus, the system is inconsistent. \textbf{No such pre-image exists.}

\textbf{Solution 9:}
The matrix is already in echelon form:
$$\begin{pmatrix} 1 & -2 & 1 & 0 & | & 3 \\ 0 & 1 & -1 & 2 & | & 4 \end{pmatrix}$$
Variables $x_1, x_2$ are basic. Variables $x_3, x_4$ are free. \\
Row 2: $x_2 - x_3 + 2x_4 = 4 \implies x_2 = 4 + x_3 - 2x_4$. \\
Row 1: $x_1 - 2x_2 + x_3 = 3$. Substitute $x_2$: \\
$x_1 - 2(4 + x_3 - 2x_4) + x_3 = 3 \implies x_1 - 8 - 2x_3 + 4x_4 + x_3 = 3 \implies x_1 - x_3 + 4x_4 = 11 \implies x_1 = 11 + x_3 - 4x_4$. \\
Parametric form:
$$\mathbf{x} = \begin{pmatrix} 11 \\ 4 \\ 0 \\ 0 \end{pmatrix} + x_3 \begin{pmatrix} 1 \\ 1 \\ 1 \\ 0 \end{pmatrix} + x_4 \begin{pmatrix} -4 \\ -2 \\ 0 \\ 1 \end{pmatrix}$$

\textbf{Solution 10 (Flaw):}
The flaw is conflating existence with uniqueness. Finding a single valid vector proves the system is consistent (a pre-image exists), but uniqueness requires checking the columns of $A$ for free variables. Since the second column is a multiple of the first (no pivot in column 2), there is a free variable, meaning there are infinitely many pre-images.

\textbf{Solution 11:}
Solve $A\mathbf{x} = \mathbf{0}$:
$$\begin{pmatrix} 3 & -6 & | & 0 \\ -2 & 4 & | & 0 \end{pmatrix}$$
$R_2 \leftarrow R_2 + \frac{2}{3}R_1 \implies \begin{pmatrix} 3 & -6 & | & 0 \\ 0 & 0 & | & 0 \end{pmatrix}$. \\
$3x_1 - 6x_2 = 0 \implies x_1 = 2x_2$. \\
$$\mathbf{x} = x_2 \begin{pmatrix} 2 \\ 1 \end{pmatrix}$$

\textbf{Solution 12:}
Solve $A\mathbf{x} = \mathbf{0}$:
$$\begin{pmatrix} 1 & 2 & -1 & | & 0 \\ 2 & 4 & -2 & | & 0 \\ -1 & -2 & 1 & | & 0 \end{pmatrix}$$
$R_2 \leftarrow R_2 - 2R_1$, $R_3 \leftarrow R_3 + R_1$:
$$\begin{pmatrix} 1 & 2 & -1 & | & 0 \\ 0 & 0 & 0 & | & 0 \\ 0 & 0 & 0 & | & 0 \end{pmatrix}$$
$x_1 + 2x_2 - x_3 = 0 \implies x_1 = -2x_2 + x_3$. \\
$$\mathbf{x} = x_2 \begin{pmatrix} -2 \\ 1 \\ 0 \end{pmatrix} + x_3 \begin{pmatrix} 1 \\ 0 \\ 1 \end{pmatrix}$$

\textbf{Solution 13:}
The matrix is in RREF:
$$\begin{pmatrix} 1 & 0 & 3 & -2 \\ 0 & 1 & -4 & 1 \\ 0 & 0 & 0 & 0 \end{pmatrix}$$
Basic: $x_1, x_2$. Free: $x_3, x_4$. \\
$x_1 = -3x_3 + 2x_4$. \\
$x_2 = 4x_3 - x_4$. \\
$$\mathbf{x} = x_3 \begin{pmatrix} -3 \\ 4 \\ 1 \\ 0 \end{pmatrix} + x_4 \begin{pmatrix} 2 \\ -1 \\ 0 \\ 1 \end{pmatrix}$$

\textbf{Solution 14:}
Solve $A\mathbf{x} = \mathbf{0}$:
$$\begin{pmatrix} 1 & 1 & 1 & 1 & | & 0 \\ 2 & 2 & 2 & 2 & | & 0 \end{pmatrix} \implies \begin{pmatrix} 1 & 1 & 1 & 1 & | & 0 \\ 0 & 0 & 0 & 0 & | & 0 \end{pmatrix}$$
$x_1 = -x_2 - x_3 - x_4$. \\
$$\mathbf{x} = x_2 \begin{pmatrix} -1 \\ 1 \\ 0 \\ 0 \end{pmatrix} + x_3 \begin{pmatrix} -1 \\ 0 \\ 1 \\ 0 \end{pmatrix} + x_4 \begin{pmatrix} -1 \\ 0 \\ 0 \\ 1 \end{pmatrix}$$

\textbf{Solution 15 (Flaw):}
The dimensions for the domain and codomain are reversed. A $5 \times 3$ matrix must multiply a $3 \times 1$ vector. Therefore, the input vector has 3 entries, meaning the domain is $\mathbb{R}^3$. The output vector will have 5 entries (one for each row), meaning the codomain is $\mathbb{R}^5$. 

\textbf{Solution 16:}
For a $7 \times 4$ matrix, the input vector $\mathbf{x}$ must be in $\mathbb{R}^4$ (number of columns) and the output will be in $\mathbb{R}^7$ (number of rows). Thus, $a = 4$ and $b = 7$.

\textbf{Solution 17:}
Since $T: \mathbb{R}^6 \to \mathbb{R}^2$, the input vectors have 6 entries and output vectors have 2 entries. By the "$m \times n$ matrix maps $\mathbb{R}^n \to \mathbb{R}^m$" rule, $n=6$ and $m=2$. The matrix $A$ must be $2 \times 6$.

\textbf{Solution 18:}
The input $\mathbf{x}$ has variables $x_1$ and $x_2$, so it is in $\mathbb{R}^2$. The output vector has 3 rows, so it is in $\mathbb{R}^3$. The domain is $\mathbb{R}^2$ and the codomain is $\mathbb{R}^3$.

\textbf{Solution 19:}
By linearity, $T(3\mathbf{u} - 2\mathbf{v}) = 3T(\mathbf{u}) - 2T(\mathbf{v})$.
$$= 3\begin{pmatrix} 2 \\ -1 \end{pmatrix} - 2\begin{pmatrix} 0 \\ 3 \end{pmatrix} = \begin{pmatrix} 6 \\ -3 \end{pmatrix} - \begin{pmatrix} 0 \\ 6 \end{pmatrix} = \begin{pmatrix} 6 \\ -9 \end{pmatrix}.$$

\textbf{Solution 20:}
By the Superposition Principle:
$$T(2\mathbf{e}_1 + 3\mathbf{e}_2 - \mathbf{e}_3) = 2T(\mathbf{e}_1) + 3T(\mathbf{e}_2) - T(\mathbf{e}_3)$$
$$= 2\begin{pmatrix} 1 \\ 4 \end{pmatrix} + 3\begin{pmatrix} -2 \\ 1 \end{pmatrix} - \begin{pmatrix} 0 \\ 5 \end{pmatrix} = \begin{pmatrix} 2 \\ 8 \end{pmatrix} + \begin{pmatrix} -6 \\ 3 \end{pmatrix} - \begin{pmatrix} 0 \\ 5 \end{pmatrix} = \begin{pmatrix} -4 \\ 6 \end{pmatrix}.$$

\textbf{Solution 21:}
$T(c_1\mathbf{v}_1 + c_2\mathbf{v}_2) = A(c_1\mathbf{v}_1 + c_2\mathbf{v}_2)$. By the distributive property of matrix multiplication: 
$$= A(c_1\mathbf{v}_1) + A(c_2\mathbf{v}_2)$$
By the scalar property of matrix multiplication:
$$= c_1(A\mathbf{v}_1) + c_2(A\mathbf{v}_2) = c_1T(\mathbf{v}_1) + c_2T(\mathbf{v}_2).$$

\textbf{Solution 22:}
Calculate left side: $\mathbf{u} + \mathbf{v} = (1+2, 1+0) = (3, 1)$. \\
$T(3, 1) = (3(3)-1, 3+4(1)) = (8, 7)$. \\
Calculate right side: $T(1, 1) = (3(1)-1, 1+4(1)) = (2, 5)$. \\
$T(2, 0) = (3(2)-0, 2+4(0)) = (6, 2)$. \\
$T(\mathbf{u}) + T(\mathbf{v}) = (2, 5) + (6, 2) = (8, 7)$. \\
Both sides equal $(8, 7)$. The property holds.

\textbf{Solution 23:}
The fastest method is the Origin Rule: $T(\mathbf{0}) = \mathbf{0}$. 
$T(0, 0) = (0 + 2, 0 - 4) = (2, -4) \neq (0, 0)$. Thus, it is not linear.

\textbf{Solution 24:}
Let $c = 2$ and $\mathbf{u} = (1, 1)$. 
Left side: $T(c\mathbf{u}) = T(2, 2) = (2^2, 2) = (4, 2)$.
Right side: $cT(\mathbf{u}) = 2 \cdot T(1, 1) = 2(1^2, 1) = 2(1, 1) = (2, 2)$.
Since $(4, 2) \neq (2, 2)$, scalar multiplication is not preserved. $T$ is not linear.

\textbf{Solution 25:}
Let $c = -1$ and $\mathbf{u} = (0, 1)$.
Left side: $T(c\mathbf{u}) = T(0, -1) = (0 - (-1), |-1|) = (1, 1)$.
Right side: $cT(\mathbf{u}) = -1 \cdot T(0, 1) = -1(0 - 1, |1|) = -1(-1, 1) = (1, -1)$.
Since $(1, 1) \neq (1, -1)$, scalar multiplication fails for negative scalars. $T$ is not linear.

\textbf{Solution 26:}
Let $\mathbf{u} = (1, 1)$ and $\mathbf{v} = (1, 1)$.
Left side: $T(\mathbf{u} + \mathbf{v}) = T(2, 2) = (2 \cdot 2, 3(2)) = (4, 6)$.
Right side: $T(\mathbf{u}) + T(\mathbf{v}) = (1 \cdot 1, 3(1)) + (1 \cdot 1, 3(1)) = (1, 3) + (1, 3) = (2, 6)$.
Since $(4, 6) \neq (2, 6)$, addition is not preserved. $T$ is not linear.

\textbf{Solution 27:}
The columns of standard matrix $A$ are simply the images of the standard basis vectors. 
$$A = [T(\mathbf{e}_1) \quad T(\mathbf{e}_2)] = \begin{pmatrix} 4 & 2 \\ -1 & 7 \end{pmatrix}.$$

\textbf{Solution 28:}
Concatenate the images of the basis vectors as columns of $A$:
$$A = [T(\mathbf{e}_1) \quad T(\mathbf{e}_2) \quad T(\mathbf{e}_3)] = \begin{pmatrix} 1 & 3 & -1 \\ 0 & 1 & -1 \\ 2 & 0 & 1 \\ -1 & 0 & 1 \end{pmatrix}.$$

\textbf{Solution 29:}
Compute $A\mathbf{x}_0$:
$$T(\mathbf{x}_0) = \begin{pmatrix} 0.85 & 0.20 \\ 0.15 & 0.80 \end{pmatrix} \begin{pmatrix} 1 \\ 0 \end{pmatrix} = \begin{pmatrix} (0.85)(1) + (0.20)(0) \\ (0.15)(1) + (0.80)(0) \end{pmatrix} = \begin{pmatrix} 0.85 \\ 0.15 \end{pmatrix}.$$
After one year, the distribution is 85\% Investment Grade and 15\% Speculative Grade.

\textbf{Solution 30:}
Compute $A\mathbf{x}_0$:
$$\mathbf{x}_1 = \begin{pmatrix} 0.5 & 0.3 & 0.1 \\ 0.3 & 0.5 & 0.2 \\ 0.2 & 0.2 & 0.7 \end{pmatrix} \begin{pmatrix} 0.4 \\ 0.4 \\ 0.2 \end{pmatrix} = \begin{pmatrix} (0.5)(0.4) + (0.3)(0.4) + (0.1)(0.2) \\ (0.3)(0.4) + (0.5)(0.4) + (0.2)(0.2) \\ (0.2)(0.4) + (0.2)(0.4) + (0.7)(0.2) \end{pmatrix}$$
$$= \begin{pmatrix} 0.20 + 0.12 + 0.02 \\ 0.12 + 0.20 + 0.04 \\ 0.08 + 0.08 + 0.14 \end{pmatrix} = \begin{pmatrix} 0.34 \\ 0.36 \\ 0.30 \end{pmatrix}.$$
Check sum: $0.34 + 0.36 + 0.30 = 1.00$. The calculation is correct.

\newpage
% ---------------------------------------------------------
% CONCEPT CHECKLIST & QUESTION MAPPING (END/REFERENCE)
% ---------------------------------------------------------
\section{Concept Checklist \& Question Mapping (Reference)}

\begin{itemize}
    \item \textbf{Type 1: Forward Mapping (Image Calculation)}
    \begin{itemize}
        \item Computing $T(\mathbf{x}) = A\mathbf{x}$ for numerical vectors. [Q1, Q2]
        \item Computing $T(\mathbf{x}) = A\mathbf{x}$ for abstract/parametric vectors. [Q3, Q4]
    \end{itemize}
    
    \item \textbf{Type 2: Inverse Mapping (Finding Pre-Images \& Uniqueness)}
    \begin{itemize}
        \item Solving $A\mathbf{x} = \mathbf{b}$ for a unique solution. [Q5, Q6]
        \item Solving $A\mathbf{x} = \mathbf{b}$ with infinite solutions (parametric form). [Q7, Q9]
        \item Solving $A\mathbf{x} = \mathbf{b}$ with no solution (inconsistent). [Q8]
        \item Diagnostic: Conflating existence with uniqueness (Flawed Problem concept). [Q10]
    \end{itemize}
    
    \item \textbf{Type 3: Kernel / Mapping to the Zero Vector}
    \begin{itemize}
        \item Solving the homogeneous equation $A\mathbf{x} = \mathbf{0}$. [Q11, Q12]
        \item Parametric vector decomposition with free variables. [Q13, Q14]
    \end{itemize}
    
    \item \textbf{Type 4: Domain and Codomain Identification}
    \begin{itemize}
        \item Matrix dimensions $m \times n$ mapping $\mathbb{R}^n \to \mathbb{R}^m$. [Q16, Q17]
        \item Reversing the rule to find matrix dimensions from mapping spaces. [Q18]
        \item Diagnostic: Dimension mismatch flaw. [Q15]
    \end{itemize}
    
    \item \textbf{Type 5: Abstract Property Application (Linearity)}
    \begin{itemize}
        \item Applying Superposition: $T(c_1\mathbf{u} + c_2\mathbf{v}) = c_1T(\mathbf{u}) + c_2T(\mathbf{v})$. [Q19, Q20, Q21]
        \item Finding images of scaled and summed vectors without matrices. [Q22]
    \end{itemize}
    
    \item \textbf{Type 6: Proving Non-Linearity}
    \begin{itemize}
        \item The "Plus a Constant" Trap / Origin Rule $T(\mathbf{0}) \neq \mathbf{0}$. [Q23, Q24]
        \item Counterexamples for scalar multiplication (absolute values, powers). [Q25, Q26]
    \end{itemize}
    
    \item \textbf{Type 7: Constructing the Standard Matrix}
    \begin{itemize}
        \item Assembling $A$ using the images of standard basis vectors $T(\mathbf{e}_1), T(\mathbf{e}_2)$. [Q27, Q28]
    \end{itemize}
    
    \item \textbf{Type 8: Real-World Applications \& Modeling}
    \begin{itemize}
        \item Markov Chains and State Vectors (Finance focus). [Q29, Q30]
    \end{itemize}
\end{itemize}

\end{document}